\documentclass[UTF8,a4paper,11pt]{ctexart}

\usepackage[a4paper,margin=2.2cm]{geometry}
\usepackage{amsmath,amssymb}
\usepackage{booktabs}
\usepackage{array}
\usepackage{longtable}
\usepackage{graphicx}
\usepackage{subcaption}
\usepackage{xcolor}
\usepackage{enumitem}
\usepackage{hyperref}
\usepackage{siunitx}
\usepackage{float}
\usepackage{caption}
\usepackage{fancyhdr}
\usepackage{lastpage}

\hypersetup{
  colorlinks=true,
  linkcolor=blue!50!black,
  urlcolor=blue!50!black,
  citecolor=blue!50!black
}

\captionsetup{font=small,labelfont=bf}
\setlist[itemize]{leftmargin=1.8em}
\setlist[enumerate]{leftmargin=2.0em}
\setlength{\parindent}{2em}
\setlength{\parskip}{0.4em}
\setlength{\headheight}{14pt}

\pagestyle{fancy}
\fancyhf{}
\lhead{mini\_grin\_rebuild 实验报告}
\rhead{\today}
\cfoot{\thepage/\pageref{LastPage}}

% ---------- Metadata ----------
\newcommand{\ReportTitle}{实验标题待填写}
\newcommand{\ReportID}{EXP-YYYYMMDD-001}
\newcommand{\AuthorName}{姓名待填写}
\newcommand{\ReportDate}{\today}
\newcommand{\ExperimentStatus}{进行中 / 已完成 / 已终止}

\newcommand{\ProblemStatement}{%
请在此处用 2--4 句话概括实验问题，例如：\\
在当前微透镜缺陷重建主线上，验证某个假设、某个数据分布变化、某个真实实验采集设置，
是否会影响 pseudo-poisson + residual 模型的稳定性、QC 门控行为和泛化能力。%
}

\newcommand{\KeyHypothesis}{%
H1: 稀疏小缺陷假设在该实验设置下仍然成立。\\
H2: \Code{standard -> test} 的差分仍足以驱动 signless / pseudo-poisson 残差学习。%
}

\newcommand{\DataSummary}{%
数据来源待填写，例如：synthetic / real pilot / hybrid。\\
若是真实实验，明确说明是否存在 `reference / standard / test` 配对。%
}

\newcommand{\PrimaryRunPaths}{%
\begin{itemize}
  \item run 路径 1：\texttt{/abs/path/to/run\_a}
  \item run 路径 2：\texttt{/abs/path/to/run\_b}
\end{itemize}
}

\newcommand{\PrimaryConfigPaths}{%
\begin{itemize}
  \item 配置 1：\texttt{/abs/path/to/config\_a.json}
  \item 配置 2：\texttt{/abs/path/to/config\_b.json}
\end{itemize}
}

\newcommand{\HardwareSummary}{%
\begin{itemize}
  \item 机器 / GPU：待填写
  \item 环境：\texttt{conda env / python version}
  \item 代码版本：git commit 待填写
\end{itemize}
}

\newcommand{\InstrumentSummary}{%
\begin{itemize}
  \item 若为仿真实验：可写 ``N/A''
  \item 若为真实实验：仪器型号、物镜、NA、波长、像素尺寸、曝光、增益、照明方式
\end{itemize}
}

% ---------- Convenience ----------
\newcommand{\TODO}[1]{\textcolor{red!70!black}{[TODO: #1]}}
\newcommand{\GOOD}[1]{\textcolor{green!40!black}{#1}}
\newcommand{\RISK}[1]{\textcolor{red!60!black}{#1}}
\newcommand{\Code}[1]{\texttt{#1}}
\newcommand{\PlaceholderGraphic}[1]{%
  \fbox{%
    \parbox[c][0.26\textheight][c]{0.9\linewidth}{%
      \centering
      \textbf{占位图}\\[0.6em]
      \texttt{\detokenize{#1}}\\[0.6em]
      请将真实实验输出图放到该路径
    }%
  }%
}

\begin{document}

\begin{titlepage}
  \centering
  {\LARGE \bfseries \ReportTitle \par}
  \vspace{1.2em}
  {\large mini\_grin\_rebuild 实验报告模板 \par}
  \vspace{2em}

  \begin{tabular}{>{\bfseries}p{0.24\textwidth} p{0.66\textwidth}}
    报告编号 & \ReportID \\
    作者 & \AuthorName \\
    日期 & \ReportDate \\
    状态 & \ExperimentStatus \\
  \end{tabular}

  \vfill

  \begin{minipage}{0.92\textwidth}
    \textbf{实验问题} \\
    \ProblemStatement
  \end{minipage}

  \vspace{1.2em}

  \begin{minipage}{0.92\textwidth}
    \textbf{关键假设} \\
    \KeyHypothesis
  \end{minipage}

  \vfill
\end{titlepage}

\tableofcontents
\newpage

\section{执行摘要}

\subsection{一句话结论}
\TODO{用 1 句话写清这次实验最重要的结论。}

\subsection{主要发现}
\begin{itemize}
  \item \TODO{发现 1}
  \item \TODO{发现 2}
  \item \TODO{发现 3}
\end{itemize}

\subsection{当前判断}
\TODO{写清这次实验对主线的意义：支持当前路线、否定某个假设、还是暴露了某个新的工程瓶颈。}

\section{背景与问题定义}

\subsection{主线位置}
建议交代这次实验处于哪条主线之下，例如：
\begin{itemize}
  \item baseline 冻结后的验证
  \item OOD / pre-real / QC / real pilot
  \item 稀疏性假设或 sign 假设压力测试
\end{itemize}

\subsection{本次实验要回答的问题}
\begin{enumerate}
  \item \TODO{问题 1}
  \item \TODO{问题 2}
  \item \TODO{问题 3}
\end{enumerate}

\subsection{验收标准}
\begin{itemize}
  \item 定量指标：\TODO{例如 global\_rmse、defect\_f1、physics\_rmse、qc\_pass\_rate}
  \item 定性指标：\TODO{例如边缘伪影、缺陷局部化、失败模式是否可解释}
  \item 决策阈值：\TODO{例如是否进入下一轮真实实验 / 是否冻结当前 baseline}
\end{itemize}

\section{实验对象与假设}

\subsection{任务定义}
\TODO{明确输入、输出、评价对象、是否依赖 reference/standard/test。}

\subsection{模型或基线假设}
建议至少覆盖：
\begin{itemize}
  \item 稀疏小缺陷假设
  \item 微透镜标准面平滑且可稳定建模
  \item `pseudo-poisson` 分支所需的符号一致性近似
  \item 真实或合成数据中 aperture / edge 的可识别性
\end{itemize}

\subsection{可能失效点}
\begin{itemize}
  \item \TODO{失效点 1}
  \item \TODO{失效点 2}
  \item \TODO{失效点 3}
\end{itemize}

\section{数据与样本说明}

\subsection{数据摘要}
\DataSummary

\subsection{数据版本与路径}
\PrimaryConfigPaths
\PrimaryRunPaths

\subsection{数据构成}
\begin{longtable}{p{0.18\textwidth} p{0.18\textwidth} p{0.16\textwidth} p{0.36\textwidth}}
\toprule
样本组 & 样本数 & 角色 & 说明 \\
\midrule
\endfirsthead
\toprule
样本组 & 样本数 & 角色 & 说明 \\
\midrule
\endhead
clean repeat & \TODO{n} & standard / repeat & \TODO{说明} \\
paired defect & \TODO{n} & standard + test & \TODO{说明} \\
reference & \TODO{n} & reference & \TODO{说明} \\
\bottomrule
\end{longtable}

\subsection{若为真实实验：采集条件}
\InstrumentSummary

\subsection{若为真实实验：元数据与配对规则}
\begin{itemize}
  \item 元数据 manifest：\Code{external\_data/templates/real\_capture\_manifest.example.json}
  \item 是否存在 \Code{reference / standard / test}：\TODO{填写}
  \item 重复采集次数：\TODO{填写}
  \item 视场限制：\TODO{单 lens / aperture 完整 / 对准条件}
\end{itemize}

\section{方法与配置}

\subsection{模型 / 基线}
\begin{itemize}
  \item 主模型：\TODO{例如 pseudo-poisson prior + UNet++ residual}
  \item 对照基线：\TODO{例如 pseudo-poisson / oracle / ablation}
\end{itemize}

\subsection{关键配置}
\begin{longtable}{p{0.34\textwidth} p{0.18\textwidth} p{0.38\textwidth}}
\toprule
参数 & 数值 & 说明 \\
\midrule
\endfirsthead
\toprule
参数 & 数值 & 说明 \\
\midrule
\endhead
\Code{scene} & \TODO{} & \TODO{} \\
\Code{grid\_size} & \TODO{} & \TODO{} \\
\Code{dx} & \TODO{} & \TODO{} \\
\Code{use\_pseudo\_poisson\_prior} & \TODO{} & \TODO{} \\
\Code{pseudo\_poisson\_prior\_as\_input} & \TODO{} & \TODO{} \\
\Code{pseudo\_poisson\_residual\_scale} & \TODO{} & \TODO{} \\
\Code{use\_phase\_inputs} & \TODO{} & \TODO{} \\
\Code{crop\_size} & \TODO{} & \TODO{} \\
\Code{model\_padding\_mode} & \TODO{} & \TODO{} \\
\bottomrule
\end{longtable}

\subsection{运行环境}
\HardwareSummary

\subsection{执行命令}
\begin{verbatim}
# 在这里记录真实执行过的命令
conda run -n PINNs bash -lc 'cd /abs/path/to/mini_grin_rebuild && ...'
\end{verbatim}

\section{指标与报告规范}

\subsection{主指标}
\begin{itemize}
  \item \TODO{例如 defect\_rmse / defect\_f1 / global\_rmse}
  \item \TODO{明确统计对象与 split}
\end{itemize}

\subsection{QC / Gating 指标}
\begin{itemize}
  \item \Code{physics\_rmse}
  \item \Code{physics\_p95\_abs}
  \item \Code{edge\_mean\_abs}
  \item \Code{edge\_p95\_abs}
  \item \Code{outside\_mean\_abs}
  \item \Code{outside\_p95\_abs}
\end{itemize}

\subsection{判定规则}
\TODO{写清 pass/fail 的依据，以及本次实验是否采用固定阈值或重新校准的 q99 阈值。}

\section{结果}

\subsection{定量结果总表}
\begin{longtable}{p{0.22\textwidth} *{5}{>{\centering\arraybackslash}p{0.13\textwidth}}}
\toprule
方法 / 实验组 & 指标 1 & 指标 2 & 指标 3 & QC 通过率 & 备注 \\
\midrule
\endfirsthead
\toprule
方法 / 实验组 & 指标 1 & 指标 2 & 指标 3 & QC 通过率 & 备注 \\
\midrule
\endhead
baseline & \TODO{} & \TODO{} & \TODO{} & \TODO{} & \TODO{} \\
method A & \TODO{} & \TODO{} & \TODO{} & \TODO{} & \TODO{} \\
method B & \TODO{} & \TODO{} & \TODO{} & \TODO{} & \TODO{} \\
\bottomrule
\end{longtable}

\subsection{关键图像}
\begin{figure}[H]
  \centering
  \begin{subfigure}[b]{0.47\textwidth}
    \centering
    \IfFileExists{figures/example_result_1.png}{%
      \includegraphics[width=\textwidth]{figures/example_result_1.png}%
    }{%
      \PlaceholderGraphic{figures/example_result_1.png}%
    }
    \caption{示例 1}
  \end{subfigure}
  \hfill
  \begin{subfigure}[b]{0.47\textwidth}
    \centering
    \IfFileExists{figures/example_result_2.png}{%
      \includegraphics[width=\textwidth]{figures/example_result_2.png}%
    }{%
      \PlaceholderGraphic{figures/example_result_2.png}%
    }
    \caption{示例 2}
  \end{subfigure}
  \caption{建议展示：标准面、预测缺陷、残差、失败案例或 QC 异常样本。}
\end{figure}

\subsection{现象解释}
\begin{itemize}
  \item \TODO{最关键现象 1}
  \item \TODO{最关键现象 2}
  \item \TODO{最关键现象 3}
\end{itemize}

\section{失败案例与风险}

\subsection{失败样本归类}
\begin{itemize}
  \item 边缘伪影主导
  \item aperture 检测失败
  \item 非稀疏表面纹理导致方法失效
  \item reference / standard / test 不可配对
  \item 真实采集漂移过大
\end{itemize}

\subsection{未解决问题}
\begin{itemize}
  \item \TODO{问题 1}
  \item \TODO{问题 2}
  \item \TODO{问题 3}
\end{itemize}

\subsection{对结论的威胁}
\begin{itemize}
  \item 样本量是否过小
  \item 是否只在单一 seed / 单一 session / 单一设备上成立
  \item 是否存在 selection bias
\end{itemize}

\section{结论与下一步}

\subsection{本次结论}
\TODO{用 3--6 句话给出收敛结论，不要重复全文。}

\subsection{是否进入下一阶段}
\begin{itemize}
  \item [ ] 可以扩大实验规模
  \item [ ] 需要先补真实数据采集
  \item [ ] 需要先修改模型假设
  \item [ ] 需要先补 QC / 数据预处理
\end{itemize}

\subsection{下一步动作}
\begin{enumerate}
  \item \TODO{动作 1}
  \item \TODO{动作 2}
  \item \TODO{动作 3}
\end{enumerate}

\appendix

\section{附录 A：实验清单}

\begin{longtable}{p{0.18\textwidth} p{0.72\textwidth}}
\toprule
项目 & 内容 \\
\midrule
\endfirsthead
\toprule
项目 & 内容 \\
\midrule
\endhead
代码版本 & \TODO{git commit} \\
数据版本 & \TODO{dataset hash / manifest id} \\
运行命令 & \TODO{命令全文} \\
关键路径 & \TODO{runs / configs / outputs} \\
环境快照 & \TODO{conda / python / torch / cuda} \\
\bottomrule
\end{longtable}

\section{附录 B：图表替换说明}

将下面这些占位图替换成你真实的输出：
\begin{itemize}
  \item \Code{figures/example\_result\_1.png}
  \item \Code{figures/example\_result\_2.png}
\end{itemize}

\end{document}
